\section{Compito 2}

\subsection{Codifica dei due algoritmi}
L'implementazione dello pseudocodice è stata fatta in C++ nella seguente maniera:
\begin{itemize}
    \item \textbf{Types.hpp}: Definisce i tipi di dati utilizzati (\texttt{CellType}, \texttt{Action}, \texttt{Position} e gli alias \texttt{Grid}, \texttt{ValueMatrix}, \texttt{PolicyMatrix});
    \item \textbf{GridMDP.hpp}: Definizione delle firme dei metodi che verranno implementati in \textbf{GridMDP.cpp};
    \item \textbf{GridMDP.cpp}: Implementazione dei metodi definiti in \textbf{GridMDP.hpp}. Le funzioni \texttt{gridValueIteration} e \texttt{constructOptimalPath} sono l'implementazione esatta dei due pseudocodici della scorsa sezione, mentre \texttt{gridValueIterationInPlace} realizza la variante descritta qui di seguito;
    \item \textbf{MemoryTracker.hpp} / \textbf{MemoryTracker.cpp}: Sostituiscono gli operatori globali \texttt{new} e \texttt{delete} per misurare il picco di memoria \textit{heap} ausiliaria allocata da ciascuna variante, dato impiegato nella sperimentazione del Compito 3;
    \item \textbf{main.cpp}: Contiene il \texttt{main} e la logica per leggere i file di input, cronometrare le due varianti e salvare i risultati;
    \item \textbf{GenerateInstance.cpp}: Programma separato, usato una tantum per produrre le 10 istanze della cartella \texttt{istanze/}. Non fa parte dell'eseguibile di risoluzione.
\end{itemize}

\noindent
La compilazione dell'eseguibile di risoluzione richiede i tre file \texttt{main.cpp}, \texttt{GridMDP.cpp} e \texttt{MemoryTracker.cpp}.

\subsection{Variante In-Place}
A differenza della versione standard (\texttt{gridValueIteration}), che calcola i nuovi valori basandosi su una copia congelata della matrice dell'iterazione precedente (\texttt{V\_old}) secondo uno schema di aggiornamento sincrono, la variante in-place \\(\texttt{gridValueIterationInPlace}) aggiorna la funzione valore in modo asincrono:
\begin{itemize}
    \item \textbf{Eliminazione della copia di matrice}: Non viene più allocata né duplicata la matrice $N \times N$ all'inizio di ogni iterazione del ciclo \texttt{while}.
    \item \textbf{Copia locale dello stato}: All'interno dei cicli di scansione, viene salvato solo il valore scalare della singola cella in esame (\texttt{v\_old = V[r][c]}), necessario per calcolare la variazione assoluta $\Delta = |v_{old} - V(s)|$.
    \item \textbf{Aggiornamento in tempo reale}: Il nuovo valore di $V(s) = \max_a Q(s, a)$ viene sovrascritto immediatamente nella matrice condivisa $V$.
\end{itemize}

\noindent
Di conseguenza, quando l'algoritmo valuta lo stato successivo $s'$ per calcolare $q = R + \gamma \cdot V(s')$, legge i dati direttamente dalla matrice $V$ in tempo reale. Se la cella adiacente $s'$ è già stata elaborata durante la scansione corrente, l'algoritmo utilizzerà il suo valore già aggiornato al ciclo attuale anziché quello dell'iterazione del ciclo \texttt{while} precedente.

\subsection{Vantaggi della Variante In-Place}
Entrambe le varianti convergono alla medesima funzione valore ottimale $V^*$, che è unica. La politica $\pi^*$ restituita può invece in linea di principio differire quando più azioni realizzano lo stesso valore massimo, poiché le due varianti visitano gli stati con valori intermedi diversi; nella sperimentazione del Compito 3 le due politiche sono comunque risultate identiche su tutte e 10 le istanze. Ciò premesso, l'approccio in-place offre due vantaggi:

\begin{enumerate}
    \item \textbf{Maggiore velocità di convergenza}:
    Nella variante standard sincrona, il valore di una cella (es. la ricompensa del Goal) può propagarsi verso i vicini al massimo di una cella per ogni iterazione del ciclo \texttt{while}; come mostrato nel Compito 1 ciò porta a $K_{\text{Standard}} = \text{ecc}(G) + 1$ iterazioni esatte. Con l'aggiornamento in-place, si innesca un effetto a catena: se l'ordine di scansione incontra una cella adiacente a una già aggiornata, recepisce immediatamente la novità. Questo permette all'informazione di propagarsi lungo più celle nella stessa iterazione, riducendo il numero totale di cicli necessari a soddisfare la condizione di arresto $\Delta < \epsilon$.
    Il guadagno dipende però dall'orientamento dei cammini ottimi rispetto all'ordine di scansione per righe: se l'informazione deve propagarsi in verso opposto alla scansione, l'effetto a catena non si innesca e le due varianti compiono lo stesso numero di iterazioni. Si ha dunque $K_{\text{In-Place}} \le K_{\text{Standard}}$, con uguaglianza nel caso sfavorevole.

    \item \textbf{Efficienza di memoria}:
    Eliminando la duplicazione continua dell'intera griglia ad ogni iterazione, la complessità spaziale \textit{ausiliaria} si riduce da $\Theta(N^2)$ a $\Theta(1)$, richiedendo soltanto la variabile temporanea $v_{old}$. Resta ovviamente invariata l'occupazione $\Theta(N^2)$ delle tabelle $V$ e $\pi$, che costituiscono l'uscita dell'algoritmo e non sono quindi memoria ausiliaria.

\end{enumerate}

\noindent
Entrambi i fenomeni vengono misurati sperimentalmente nel Compito 3.
